\documentclass[10pt,a4paper]{amsart}
\usepackage[margin=19mm]{geometry}
\usepackage{amsmath,amssymb,mathtools}
\usepackage[scr=rsfs,cal=euler]{mathalfa}
\usepackage{xfrac}
\usepackage[unicode,hidelinks,bookmarks=false]{hyperref}
\newcommand{\ie}{i.e.\ }
\newcommand{\cf}{cf.\ }
\newcommand{\resp}{respectively\ }
\newcommand{\Subsection}{\subsection}
\providecommand{\nobreakdash}{\nobreak\hbox{-}}
\setlength{\emergencystretch}{2em}
\multlinegap0pt
\usepackage{mathtools}
\usepackage{xcolor}
\newtheorem{proposition}{Proposition}
\newcommand{\apx}{\mathrm{apx}}
\newcommand{\iapx}{I_{\mathrm{apx}}}
\providecommand{\paragraph}{}
\title{State-Dependent Lyapunov Analysis of Rank-1 Matrix Factorization}
\author{Jaehong Moon}
\date{Source: arXiv:2604.26993v3 (CC BY 4.0)}
\begin{document}
\maketitle

\noindent\emph{Source and licence.} State-Dependent Lyapunov Analysis of Rank-1 Matrix Factorization. Version arXiv:2604.26993v3 (CC BY 4.0), distributed by arXiv under the Creative Commons Attribution 4.0 International licence, http://creativecommons.org/licenses/by/4.0/, as recorded in the arXiv licence metadata for this exact version and shown on the abstract page https://arxiv.org/abs/2604.26993v3. Full licence notice: \texttt{licenses/arxiv-2604.26993-CC-BY-4.0.txt}.\par\smallskip

\noindent\textit{Editorial scope.} Counted unit: the article's proposition prop:boundary\_inward\_rank1\_approx (source0.tex lines 720--728). The article's complete original proof of it is delivered with it (source0.tex lines 2048--2218, one \texttt{proof} environment of 171 lines, which the article places away from the statement). No mathematics is written, repaired or re-derived: the statement and the proof are the article's own lines, in the article's own notation, and no retained line is substituted or re-flowed.  No further source-local context is needed: every cross-reference of every retained line resolves inside the two delivered spans.  Nothing was omitted from any retained span, so the package carries no omission record.  No retained line contains a citation, so the wrapper carries no bibliography and no bibliography entry.  No retained line contains a figure environment, an image-inclusion command or drawing code, so no image is delivered and the package declares no asset dependency.  Editorial formatting only, in addition: an AMS \texttt{amsart} preamble that restates the article's own declarations and macros used by the retained lines, namely the environments \texttt{proposition} on separate counters, together with the standard packages those lines need.  The retained environments print in this document as Proposition 1 in order of appearance, whereas the article prints the same items with the numbering of its own full document.  The complete native source of the article as distributed in the arXiv e-print for this exact version is bundled unchanged as \texttt{sources/199439/source0.tex}.
\begin{proposition}\label{prop:boundary_inward_rank1_approx}
Let $0 < \eta < \delta < 2$ and suppose $q_\eta(\delta) < 0$.
If $\iapx(\delta;\, A_t, B_t) = 0$ and $(a_t, b_t, u_t, v_t) \notin \mathcal{S}$,
then
\begin{align}
    \iapx(\delta;\, A_{t+1}, B_{t+1}) = R_t^\apx(\delta) < 0,
\end{align}
and hence the next iterate lies strictly inside the sublevel set $\{\iapx(\delta) \le 0\}$.
\end{proposition}

\begin{proof}[Proof of Proposition~\ref{prop:boundary_inward_rank1_approx}]
Fix $0<\eta<\delta<2$ with
\begin{align}
    q_\eta(\delta):=\eta\delta^2-4\delta+4\eta<0,
\end{align}
and consider a state $(a_t,b_t,u_t,v_t)$ such that
\begin{align*}
    \iapx(\delta;\,A_t,B_t)=0.
\end{align*}
If this point lies in $\mathcal S_\apx$, then the gradient step fixes it, and the
quotient--remainder identity gives $R_t^\apx(\delta)=0$. We prove that
$R_t^\apx(\delta)<0$ at every boundary point outside $\mathcal S_\apx$.

Using
\begin{align}
    D_t
    =D_t^S+\|a_t\|^2v_t^2+\|b_t\|^2u_t^2+u_t^2v_t^2,
\end{align}
we may rewrite the remainder as
\begin{align}
    R_t^\apx(\delta)
    &=-\eta(\delta-\eta)r_t(\delta)
      -\eta^2\delta^2D_t^N
      +\eta q_\eta(\delta)D_t^S,
      \label{eq:R_t_apx_regroup}\\
    r_t(\delta)
    &:=(4-\delta^2)L_t^2
      +4\bigl(\|a_t\|^2v_t^2+\|b_t\|^2u_t^2+u_t^2v_t^2\bigr)
      -\delta N_t.
      \label{eq:r_t_apx_regroup}
\end{align}
Since $D_t^N,D_t^S\ge0$ and $q_\eta(\delta)<0$, the last two terms in
Eq.~\eqref{eq:R_t_apx_regroup} are nonpositive, and their sum vanishes only if
$D_t^N=D_t^S=0$. It remains to prove $r_t(\delta)\ge0$ on
the boundary $\{\iapx(\delta)=0\}$ and to track its equality case.

\paragraph{Reduction to $u_tv_t=0$.}
Fix $(a_t,b_t)$ and $N_t=u_t^2+v_t^2$. As before, the constraint
$\iapx(\delta;\,A_t,B_t)=0$ is unchanged when $(u_t,v_t)$ varies with
$N_t$ fixed. Using
$v_t^2=N_t-u_t^2$, the off-signal expression in $r_t$ becomes
\begin{align}
    \|a_t\|^2(N_t-u_t^2)
    +\|b_t\|^2u_t^2
    +u_t^2(N_t-u_t^2).
\end{align}
Viewed as a function of $u_t^2\in[0,N_t]$, this is a concave quadratic, so its
minimum is attained at $u_t^2=0$ or $u_t^2=N_t$. Thus the minimum of $r_t$ at
fixed $(a_t,b_t,N_t)$ occurs when $u_tv_t=0$. By the exchange symmetry, it
suffices to analyze $u_t=0$, in which case
\begin{align}
    r_t(\delta)
    =(4-\delta^2)L_t^2+(4\|a_t\|^2-\delta)v_t^2.
    \label{eq:r_reduced_apx}
\end{align}

\paragraph{The case $4\|a_t\|^2-\delta>0$.}
Both terms in Eq.~\eqref{eq:r_reduced_apx} are nonnegative. If
$R_t^\apx(\delta)=0$, all three nonpositive terms in
Eq.~\eqref{eq:R_t_apx_regroup} must vanish. In particular,
$r_t(\delta)=D_t^S=0$, so $L_t=v_t=0$ and $a_t,b_t$ are linearly dependent.
Together with $u_t=0$, these conditions place the point in $\mathcal M_\apx$.

\paragraph{The case $4\|a_t\|^2-\delta=0$.}
Here $r_t(\delta)=(4-\delta^2)L_t^2\ge0$. Equality in
Eq.~\eqref{eq:R_t_apx_regroup} would force $L_t=D_t^S=0$. Thus
$\langle a_t,b_t\rangle=1$, and equality in the Cauchy--Schwarz inequality,
together with $\|a_t\|^2=\delta/4$, gives $\|b_t\|^2=4/\delta$.
The identity $\iapx(\delta;\,A_t,B_t)=0$ would then imply
\begin{align}
    v_t^2
    =\frac{4-\delta^2}{\delta}
     -\frac{\delta}{4}-\frac{4}{\delta}+\delta
    =-\frac{\delta}{4},
\end{align}
a contradiction. Hence equality in the remainder cannot occur in this case.

\paragraph{The case $4\|a_t\|^2-\delta<0$ and $\|a_t\|>0$.}
Set
\begin{align}
    \gamma:=\frac{\langle a_t,b_t\rangle}{\|a_t\|^2}.
\end{align}
The Cauchy--Schwarz inequality gives
$\|b_t\|^2/\|a_t\|^2\ge\gamma^2$. Together with
$\iapx(\delta;\,A_t,B_t)=0$, this gives
\begin{align}
    v_t^2
    &=\frac{4-\delta^2}{\delta}
      -\|a_t\|^2
       \left(1+\frac{\|b_t\|^2}{\|a_t\|^2}-\delta\gamma\right)
      \notag\\
    &\le \frac{4-\delta^2}{\delta}
       -\|a_t\|^2(1+\gamma^2-\delta\gamma).
    \label{eq:v_sq_upper_apx}
\end{align}
The denominator below is positive because
\begin{align}
    1+\gamma^2-\delta\gamma
    =\left(\gamma-\frac{\delta}{2}\right)^2
     +1-\frac{\delta^2}{4}
    >0.
\end{align}
Together with $v_t^2\ge0$, this yields
\begin{align}
    0<\|a_t\|^2\le
    \frac{4-\delta^2}
         {\delta(1+\gamma^2-\delta\gamma)}.
    \label{eq:a_norm_upper_apx}
\end{align}
Because $4\|a_t\|^2-\delta<0$, multiplying the upper bound
Eq.~\eqref{eq:v_sq_upper_apx} by this coefficient reverses the inequality.
Using $L_t=1-\|a_t\|^2\gamma$ in Eq.~\eqref{eq:r_reduced_apx}, we obtain
\begin{align}
    r_t(\delta)
    \ge \|a_t\|^2\left[
        \frac{(\delta\gamma^2-4\gamma+\delta)^2}
             {\delta(1+\gamma^2-\delta\gamma)}
        +\left(
            \frac{4-\delta^2}
                 {\delta(1+\gamma^2-\delta\gamma)}-\|a_t\|^2
         \right)(\delta\gamma-2)^2
      \right]
    \ge0,
    \label{eq:r_apx_square_bound}
\end{align}
where the last inequality follows from
Eq.~\eqref{eq:a_norm_upper_apx}.

Suppose now that $R_t^\apx(\delta)=0$. Then
Eq.~\eqref{eq:R_t_apx_regroup} forces $r_t(\delta)=D_t^S=0$.
Equality in Eq.~\eqref{eq:r_apx_square_bound} first gives
$\delta\gamma^2-4\gamma+\delta=0$. As in the factorization case, this cannot
hold together with $\delta\gamma-2=0$, since the two equalities would force
$\delta=2$. Thus equality must also hold in Eq.~\eqref{eq:a_norm_upper_apx}.
The identity
\begin{align}
    \delta(1+\gamma^2-\delta\gamma)
    -(4-\delta^2)\gamma
    =\delta\gamma^2-4\gamma+\delta
\end{align}
and $\delta\gamma^2-4\gamma+\delta=0$ give
$\delta(1+\gamma^2-\delta\gamma)=(4-\delta^2)\gamma$. Combining this with
equality in Eq.~\eqref{eq:a_norm_upper_apx} yields
$\|a_t\|^2=1/\gamma$, and hence $\|a_t\|^2\gamma=1$. Since $D_t^S=0$,
equality holds in Cauchy--Schwarz, so Eq.~\eqref{eq:v_sq_upper_apx} is an
equality; together with
Eq.~\eqref{eq:a_norm_upper_apx}, this gives $v_t=0$. Thus $u_t=v_t=0$,
$D_t^S=0$, and $\langle a_t,b_t\rangle=1$, which places the point in
$\mathcal M_\apx$.

\paragraph{The case $\|a_t\|=0$.}
Here $L_t=1$, and $\iapx(\delta;\,A_t,B_t)=0$ gives
\begin{align}
    \delta(\|b_t\|^2+v_t^2)+\delta^2-4=0.
\end{align}
Consequently,
\begin{align}
    r_t(\delta)
    =4-\delta^2-\delta v_t^2
    =\delta\|b_t\|^2\ge0.
\end{align}
Equality holds only if $b_t=0$, in which case
$(a_t,b_t,u_t,v_t)=(0,0,0,v_t)\in\mathcal S_\apx$.

The other endpoint, $v_t=0$, follows by the exchange symmetry. We have proved
$r_t(\delta)\ge0$ on $\{\iapx(\delta)=0\}$. All three terms in
Eq.~\eqref{eq:R_t_apx_regroup} are therefore nonpositive. If the remainder
vanishes, then $D_t^N=D_t^S=0$, so one of the two endpoint analyses applies,
and the case analysis above places the point in $\mathcal S_\apx$. Hence every
boundary point outside $\mathcal S_\apx$ satisfies $R_t^\apx(\delta)<0$, as claimed.
\end{proof}

\end{document}
